% Chapter Template

\chapter{Introduction}\doublespacing % Main chapter title

\label{Chapter1} % Change X to a consecutive number; for referencing this chapter elsewhere, use \ref{ChapterX}

\lhead{Chapter I. \emph{Introduction}} % Change X to a consecutive number; this is for the header on each page - perhaps a shortened title

%----------------------------------------------------------------------------------------
%	SECTION 1
%----------------------------------------------------------------------------------------
\section{Introduction}

\begin{frame}

Malware continues to evolve rapidly, posing serious threats to modern computer systems. 
Traditional signature-based detection methods often fail to detect new or obfuscated malware. 
Therefore, malware analysis has become essential for understanding malicious behavior and attack patterns. 
However, the raw data generated during malware analysis is often complex and difficult to interpret. 
\textbf{Malviz} addresses this challenge by transforming malware behavior into intuitive visual patterns, making analysis faster and more accessible for security analysts.

\end{frame}

\section{Problem Statement}

\begin{frame}

Malware is becoming increasingly sophisticated and difficult to detect using traditional security methods. 
Signature-based detection techniques frequently fail to identify new, unknown, or polymorphic malware. 
Malicious programs often perform hidden system-level activities such as privilege escalation, suspicious system calls, and abnormal process behavior. 
Manually analyzing these behaviors from raw logs and analysis reports can be complex, time-consuming, and inefficient for security analysts.

\end{frame}


 \section{Need of the Study}

With the rapid growth of digital infrastructure, malware attacks have become more frequent and sophisticated. 
Traditional malware detection techniques mainly rely on signature-based approaches, which are often ineffective against new or obfuscated malware variants. 
Security analysts usually depend on large volumes of logs, system calls, and behavioral reports generated during malware analysis. 
However, interpreting this raw data can be complex, time-consuming, and prone to human error.

There is a growing need for tools that can simplify malware analysis and help analysts quickly understand malicious behavior. 
Visual representation of malware activity can make patterns easier to identify and improve the efficiency of security investigations. 
Therefore, this study focuses on developing \textbf{Malviz}, a system that converts malware behavior into intuitive visual patterns, enabling faster and more effective malware analysis.


\section{Study Objective}

The primary objective of this study is to design and develop a system that visualizes malware behavior to simplify the analysis process for security researchers and analysts.

\begin{itemize}
    \item To analyze malware behavior based on system-level activities such as system calls, process actions, and privilege changes.
    \item To transform raw malware analysis data into meaningful visual representations.
    \item To help analysts quickly identify suspicious patterns and abnormal behavior.
    \item To provide an intuitive and accessible platform for malware behavior visualization.
\end{itemize}


\section{Dissertation Organization}

This dissertation consists of six chapters, including the present chapter (Chapter \ref{Chapter1}), which introduces the research topic. 
This chapter describes the need for the study, the problem statement, and the objectives of the research. 

Chapter \ref{Chapter2} presents a review of existing literature related to malware analysis techniques, malware visualization approaches, and cybersecurity tools used for behavior analysis. 

Chapter \ref{Chapter3} explains the methodology adopted for the development of the \textbf{Malviz} system, including data collection, malware behavior extraction, and visualization techniques. 

Chapter \ref{Chapter4} describes the implementation details of the proposed system, including system architecture, tools used, and data processing techniques. 

Chapter \ref{Chapter5} presents the results obtained from testing the system and discusses how visualization helps in understanding malware behavior. 

Finally, Chapter \ref{Chapter6} summarizes the conclusions of the research and highlights the contributions and possible future improvements.


\section{Summary}

This chapter introduced the background and motivation for the study and explained the need for improved malware analysis methods. 
The objectives of the research were presented, focusing on the development of the \textbf{Malviz} system for malware behavior visualization. 
In addition, the overall structure of the dissertation was outlined to provide a clear overview of the research organization.
 